% English translation, made from our own transcription (original.tex), not from any existing
% modern translation. All mathematics, equation numbers (\tag), \origpage markers, and figures
% are preserved unchanged; only the prose is translated. "Schediasma" (Gk. improvised essay) is
% rendered "Sketch" in the translation (the transcription keeps the original term); "arc." (the
% author's abbreviation for arco/arc) is kept. The eq. (2) misprint and the "zz" for z^2 are
% reproduced from our transcription, not corrected.
\documentclass{article}
\usepackage{readmasters}
\begin{document}

\origpage{293}
\section*{XXXII. Method for measuring the lemniscate}
\subsection*{Sketch I (*)}
\emph{(*) Giornale de' letterati d'Italia, vol.~XXIX, p.~258.}

The two eminent geometers Signor Giacomo and Signor Giovanni, the brothers Bernoulli, have
made the lemniscate famous, using its arcs to construct the paracentric isochrone, as may be
seen in the Acta of Leipzig of the year 1694. It is evident that, by measuring the lemniscate
by means of some other curve simpler than it, one obtains a more perfect construction not only
of the paracentric isochrone, but also of the other infinite curves which, in order to be
constructed, may depend on the lemniscate; and therefore I flatter myself that the measures of
this curve discovered by me, which I shall set forth successively in two sketches, will not
displease those versed in the matter.

\emph{Assumptions known to those versed in the infinitesimal calculus.} --- Let $CQACFC$
(fig.~24) be the lemniscate, whose semi-axis $CA = a$; it is known that, taking the origin of
the abscissae $(x)$ at the centre $C$, and calling $(y)$ the ordinates normal to the axis, the
nature of the lemniscate is expressed by this equation $x^{2} + y^{2} = a\sqrt{x^{2} - y^{2}}$.

It is also known that if one calls $z$ the indeterminate chord $CQ = \sqrt{x^{2} + y^{2}}$, one
has the direct arc $CQ = \displaystyle\int \frac{a^{2}\,dz}{\sqrt{a^{4} - z^{4}}}$, and the
inverse arc
\[ QA = \mathrm{arc.}\,CA - \mathrm{arc.}\,CQ = \int \frac{a^{2}\,dz}{\sqrt{a^{4} - z^{4}}}. \]

Let $ADFNA$ (fig.~24) be the ellipse, whose minor semi-axis $CF = a$, and major semi-axis
$CD = a\sqrt{2}$; call $z$ the indeterminate abscissa $CH$, which has its origin at the centre
$C$ of the ellipse and is equal to the chord $CQ$ of the lemniscate, and let the ordinate $HI$
be drawn parallel to the major axis; it is already known that the direct arc $DI$ of this
ellipse has for its expression
\[ \int dz\, \frac{\sqrt{a^{2} + z^{2}}}{\sqrt{a^{2} - z^{2}}}, \]
\origpage{294}
and that the inverse arc
\[ IF = \mathrm{arc.}\,DF - \mathrm{arc.}\,DI = \int -dz\, \frac{\sqrt{a^{2} + z^{2}}}{\sqrt{a^{2} - z^{2}}}. \]

\rmfigure{figures/fig-24.png}{Fig.~24.}{Fig. 24: the lemniscate CQACFC and its circumscribed ellipse ADFNA, from the original plate (Tav. II, Opere Matematiche vol. II).}

Let there be, finally, the equilateral hyperbola $LMP$ (fig.~25), whose semi-axis $SM = a$. Let
$t$ be named the indeterminate ordinate $SO$, which, starting from the centre $S$, cuts the arc
$MO$; it is known that this same arc is expressed thus
\[ \int \frac{t^{2}\,dt}{\sqrt{t^{4} - a^{4}}}. \]

\rmfigure{figures/fig-25.png}{Fig.~25.}{Fig. 25: the equilateral hyperbola LMP with ordinate SO and arc MO, from the original plate (Tav. III, Opere Matematiche vol. II).}

\textbf{Theorem I.} --- Let the two equations (1) and (2) written below be given. I say that,
the first of them being posited, the other also holds

\[ t = a\, \frac{\sqrt{a^{2} + z^{2}}}{\sqrt{a^{2} - z^{2}}}. \tag{1} \]

\[ \int \frac{a^{2}\,dz}{\sqrt{a^{4} - z^{4}}} = \int dz\, \frac{\sqrt{a^{2} + zz}}{\sqrt{a^{2} - az}} + \int \frac{t^{2}\,dt}{\sqrt{t^{4} - a^{4}}} - \frac{zt}{a}. \tag{2} \]

The truth of this theorem is demonstrated by differentiating, and substituting in place of $t$
and $dt$ their values in $z$ and $dz$ drawn from equation (1).

\textbf{Corollary} (fig.~24 and 25). --- If in the lemniscate the chord $CQ = z$, and in the
ellipse the abscissa $CH$ is likewise $= z$, and in the equilateral hyperbola $LMP$ the central
ordinate $SO = t$; assigning to $t$ its value expressed in equation (1), and putting in
equation (2) the arcs of the curves instead of their expressions already declared in the
assumptions above, one obtains

\[ \mathrm{Arc.}\,CQ = \mathrm{arc.}\,DI + \mathrm{arc.}\,MO = \frac{zt}{a}. \tag{3} \]

\textbf{Scholium I.} --- Supposing $z = 0$, the direct arcs $CQ$ of the lemniscate and $DI$ of
the ellipse are equal to zero; the hyperbolic arc $MO$ is likewise null, since by virtue of
equation (1) the vanishing of $z$ makes $SO\,(t) = a = SM$; the rectilinear expression is also
annihilated, and all this is a sure indication that equation (3) is complete and needs no
addition or subtraction of any constant quantity. But putting $z = a$, then the direct arc of
the lemniscate becomes equal to the arc $CQA$, and the elliptic arc $DI$ to the arc $DIF$; but
by equation (1) the central ordinate $SO\,(t)$ of the hyperbola becomes in this case infinite,
rendering the hyperbolic arc $MO$ infinite as well, and therefore it seems impossible to arrive
by this route at knowing the
\origpage{295}
value of the quadrant of the lemniscate expressed in finite quantities, and consequently the
measure of this curve cannot yet be said to be perfectly discovered.

\textbf{Theorem II.} --- Let the two equations (4) and (5) written below be given; I say that,
the first of them being posited, the other also holds

\[ r = \frac{a^{2}}{z}. \tag{4} \]

\[ \int -\frac{a^{2}\,dz}{\sqrt{a^{4} - z^{4}}} = \int -dz\, \frac{\sqrt{a^{2} + z^{2}}}{\sqrt{a^{2} - z^{2}}} + \int \frac{r^{2}\,dr}{\sqrt{r^{4} - a^{4}}} - \frac{1}{z}\sqrt{a^{4} - z^{4}}. \tag{5} \]

The differentiation and the substitution of the values of $r$ and $dr$ in $z$ and $dz$ deduced
from equation (4) demonstrate this theorem.

\textbf{Corollary I} (fig.~24 and 25). --- Let there be made, as above, in the lemniscate the
chord $CQ = z$, in the ellipse the abscissa $CH$ likewise $= z$, and in the hyperbola the
central ordinate $SL = r$, attributing to $r$ its value noted in equation (4), and substituting
in equation (5) the arcs of the curves in place of their expressions, as was done in the
corollary of Theorem I, one finds again

\[ \mathrm{Arc.}\,QA = \mathrm{arc.}\,IF + \mathrm{arc.}\,ML - \frac{1}{z}\sqrt{a^{4} - z^{4}}. \tag{6} \]

\textbf{Scholium II.} --- The supposition $z = a$ makes null the inverse arcs $QA$ of the
lemniscate and $IF$ of the ellipse; it also annuls the rectilinear expression and the
hyperbolic arc $ML$, seeing that, by force of equation (4), the central ordinate of the
hyperbola $SL\,(r)$ becomes equal to the semi-axis $SM\,(a)$; hence equation (5) is complete.
But the supposition $z = 0$ prolongs to infinity the central ordinate of the hyperbola, and its
arc, at the same time as it makes the inverse arcs of the lemniscate and of the ellipse equal
respectively to the quadrants $ACQ$, $FID$; and here the same difficulty reappears that was set
out at the end of Scholium I; but this difficulty will be overcome by the corollary that
follows. Meanwhile let it be noted: first, that by drawing in the hyperbola the $OP$ parallel to
the second axis, the arc $LMP$ is equal to the two arcs $MO$, $ML$. Secondly, that assigning to
$t$ its value expressed in equation (1), this complex quantity
$\dfrac{zt}{a} + \dfrac{1}{z}\sqrt{a^{4} - z^{4}}$ is equivalent to this other, simpler quantity
$\dfrac{at}{z}$, as will be verified by the calculation.
\origpage{296}
\textbf{Corollary II} (fig.~24 and 25). --- Adding the two equations (3) and (6), and making use
of the two remarks noted at the end of the preceding scholium, one discovers
$\mathrm{arc.}\,CQA = \mathrm{arc.}\,DIF + \mathrm{arc.}\,LMP - \dfrac{at}{z}$; and finally, by
multiplying this last equation by 4, one finds again that the entire perimeter of the lemniscate
is equal to the entire perimeter of the circumscribed ellipse $ADFNA$, plus four times the arc
$LMP$ of the equilateral hyperbola minus $\dfrac{4at}{z}$ --- which is a new and excellent
property of these curves.

\textbf{Theorem III.} --- Let the two equations (7) and (8) written below be given; I say that,
the first of them being posited, the other also holds

\[ u = a\, \frac{\sqrt{a^{2} - z^{2}}}{\sqrt{a^{2} + z^{2}}}. \tag{7} \]

\[ \int \frac{a^{2}\,dz}{\sqrt{a^{4} - z^{4}}} = \int -\frac{a^{2}\,du}{\sqrt{a^{4} - u^{4}}}. \tag{8} \]

This theorem is demonstrated, like the other two that precede it, by differentiating and making
the due substitutions drawn from equation (7).

\textbf{Corollary I} (fig.~24). --- Let there be taken in the lemniscate the chord $CQ = z$, and
the other chord $CE = u$; let there be attributed to $u$ its value expressed in equation (7);
let the direct and inverse arcs of the curve be substituted in equation (8) in place of their
expressions, and one will find $\mathrm{arc.}\,CQ = \mathrm{arc.}\,EA$.

\textbf{Scholium III} (fig.~24). --- If $z = 0$ the direct arc $CQ$ is null, and by virtue of
equation (7) the vanishing of $z$ makes $CE\,(u) = a$, so that the inverse arc $EA$ is null as
well; which makes known that the equation noted at the end of the preceding corollary is
complete. This notable property of the lemniscate was discovered by me, and set forth in a
different manner in my method for rectifying the difference of the arcs in infinite kinds of
non-rectifiable parabolas.

\textbf{Corollary II} (fig.~24). --- Finding, as I have taught above to do, the inverse arc $EA$
equal to the direct arc $CQ$, this same arc $CQ$ may be measured by means of Corollary I of
Theorem II; and conversely, finding the direct arc $CQ$ equal to the inverse arc $EA$, this same
arc $EA$ may be measured by means of the corollary of Theorem I.
\origpage{297}
\textbf{Corollary III.} --- \textbf{Problem.} --- To bisect the quadrant of the lemniscate
(fig.~24).

If the chord $CQ\,(z)$ of the lemniscate is equal to the other chord $CE\,(u)$ of it, and if at
the same time $u$ has its value noted in equation (7), then
$z = u = a\, \dfrac{\sqrt{a^{2} - z^{2}}}{\sqrt{a^{2} + z^{2}}}$, whence there results, the due
operations being performed,
\[ z = u = \sqrt{a^{2}\sqrt{2} - a^{2}} \]
and therefore, attributing to $z$ this last value, the two points of the lemniscate $Q$ and $E$
will coincide in $B$, in such a way that the direct arc $CB$ will be equal to the inverse arc
$BA$, and the chord $CB = u = z$ will bisect the quadrant of the lemniscate. Which was to be
found.

\end{document}
